
%(BEGIN_QUESTION)
% Copyright 2010, Tony R. Kuphaldt, released under the Creative Commons Attribution License (v 1.0)
% This means you may do almost anything with this work of mine, so long as you give me proper credit

\noindent
{\bf Labøvelse}

\vskip 5pt

Lagets oppgave er å bygge en komplett, fungerende prosess som styres av en nettverkstilkoblet digital regulator eller et distribuert kontrollsystem.  Alternativer inkluderer enkeltløkke-regulatorer koblet til en arbeidsstasjon via nettverkskabling, nettverkstilkoblede PLC-er som implementerer PID-regulering, DDC-enheter (byggautomasjon), eller et DCS.  Når dere kjører systemet, må dere demonstrere operatørstyring (f.eks. endring av settpunkt, auto/manuell-modus, osv.) fjernstyrt via en nettverkstilkoblet PC til prosessregulatoren.

Hvert instrument i sløyfen skal merkes med et korrekt tag-navn (f.eks. "LT-59" for en nivåtransmitter), og alle instrumenter i hver sløyfe skal ha samme sløyfenummer.  Skriv på biter av maskeringstape for å lage enkle merker for alle instrumenter og signallinjer.

Etter at hele reguleringssløyfen er bygget, skal dere "tune" P-, I- og D-innstillingene på regulatoren slik at den viser god reguleringsoppførsel (definert her som kvartbølge-demping eller bedre).

\vskip 10pt

\underbar{Tabell for måloppnåelse:}

% No blank lines allowed between lines of an \halign structure!
% I use comments (%) instead, so that TeX doesn't choke.

$$\vbox{\offinterlineskip
\halign{\strut
\vrule \quad\hfil # \ \hfil & 
\vrule \quad\hfil # \ \hfil & 
\vrule \quad\hfil # \ \hfil & 
\vrule \quad\hfil # \ \hfil & 
\vrule \quad\hfil # \ \hfil & 
\vrule \quad\hfil # \ \hfil & 
\vrule \quad\hfil # \ \hfil \vrule \cr
\noalign{\hrule}
%
% First row
{\bf Prestasjonsmål} & {\bf Vurdering} & {\bf 1} & {\bf 2} & {\bf 3} & {\bf 4} & {\bf Team} \cr
%
\noalign{\hrule}
%
% Another row
Velg prosess som skal bygges & mastery & -- & -- & -- & -- & \cr
%
\noalign{\hrule}
%
% Another row
Finn og naviger i manual for regulator & mastery & -- & -- & -- & -- & \cr
%
\noalign{\hrule}
%
% Another row
Sløyfediagram og prosessinspeksjon & mastery & & & & & -- -- -- -- \cr
%
\noalign{\hrule}
%
% Another row
Prosessen viser god reguleringsoppførsel & mastery & -- & -- & -- & -- &  \cr
%
\noalign{\hrule}
%
% Another row
Labspørsmål: Valg/testing & proportional &  &  &  &  & -- -- -- -- \cr
%
\noalign{\hrule}
%
% Another row
Labspørsmål: Igangkjøring & proportional &  &  &  &  & -- -- -- -- \cr
%
\noalign{\hrule}
%
% Another row
Labspørsmål: Hoderegning & proportional &  &  &  &  & -- -- -- -- \cr
%
\noalign{\hrule}
%
% Another row
Labspørsmål: Diagnostikk & proportional &  &  &  &  & -- -- -- -- \cr
%
\noalign{\hrule}
} % End of \halign 
}$$ % End of \vbox

Hver student vil bli bedt om å svare riktig på et "labspørsmål" fra hver av de fire kategoriene (eksempler vist på neste side).  Disse labspørsmålene fungerer som en guide til kunnskap og ferdigheter alle i teamet bør lære etter hvert som dere jobber med labøvelsen.  Instruktøren kan spørre studenter om disse spørsmålene når som helst før labøvelsen er fullført.

\vfil \eject

\noindent
\underbar{Labspørsmål} 

\begin{itemize}
\item{} {\bf Valg og innledende testing}
\item{} Identifiser alle innganger og utganger på feltinstrumentene (transmitter og FCE)
\item{} Forklar betydningen av de ulike klassifiseringene som er spesifisert på instrumentets merkeskilt
\item{} Finn i produsentdokumentasjonen hvor signalledninger skal kobles til feltinstrumentet (transmitter eller FCE)
\item{} Forklar hvilke typer testutstyr som ble brukt for å verifisere funksjonen til feltinstrumentet (transmitter eller FCE)
\item{} Forklar hvordan du kan utføre enkle tester av instrumentfunksjon ved å bruke enkelt testutstyr (multimeter, luftpumper, trykkmålere, motstander, batterier, osv.)
\end{itemize}

\filbreak

\begin{itemize}
\item{} {\bf Igangkjøring og dokumentasjon}
\item{} Vis hvordan du isolerer potensielt farlig energi i systemet ({\it lock-out, tag-out}) og også hvordan du sikkert verifiserer at energien er isolert før arbeidet starter
\item{} Vis hvordan du bruker en sløyfekalibrator til å måle/gi/simulere signalstrøm
\item{} Vis riktig valg og bruk av fastnøkler på muttere, bolter og/eller rørkoblinger
\item{} Identifiser nettverksforbindelsen(e) brukt for å få fjernaksess til regulatorvariabler fra en PC (f.eks. IP-adresser, Ethernet-forbindelser, Modbus-nettverk og adresser, osv.)
\item{} Identifiser flere steder (med henvisning til sløyfediagram) der du kan måle ulike 4-20 mA instrument-signaler i systemet
\item{} Identifiser flere steder (med henvisning til sløyfediagram) der du kan koble til HART-kommunikator i systemet
\end{itemize}

\filbreak

\begin{itemize}
\item{} {\bf Hoderegning} (ingen kalkulator tillatt!)
\item{} Bestem tillatt kalibreringsfeil for instrumentet (f.eks. +/- 0.5\% for et instrument som er ranged 200 til 500 grader)
\item{} Konverter et 4-20 mA-signal til prosent av spenn (f.eks. 13 mA = \underbar{\hskip 20pt}\%)
\item{} Konverter prosent av spenn til en 4-20 mA-signalkverdi (f.eks. 70\% = \underbar{\hskip 20pt} mA)
\item{} Konverter et 3-15 PSI-signal til prosent av spenn (f.eks. 11 PSI = \underbar{\hskip 20pt}\%)
\item{} Konverter prosent av spenn til en 3-15 PSI-signalkverdi (f.eks. 40\% = \underbar{\hskip 20pt} PSI)
\end{itemize}

\filbreak

\begin{itemize}
\item{} {\bf Diagnostikk}
\item{} Gitt en bestemt komponent- eller kablingsfeil ({\it instruktøren spesifiserer type og plassering}), hvilke symptomer vil sløyfen vise og hvorfor?
\item{} Forklar hvorfor det å bryte en 4-20 mA-sløyfe kan forårsake alvorlige problemer i en faktisk instrumentering!
\item{} Forklar hva som vil skje (og hvorfor) i reguleringssløyfen hvis transmitteren plutselig feiler med et lavt (4 mA) signal.  Anta at regulatoren står i automatisk modus når dette skjer.
\item{} Forklar hva som vil skje (og hvorfor) i reguleringssløyfen hvis transmitteren plutselig feiler med et høyt (20 mA) signal.  Anta at regulatoren står i automatisk modus når dette skjer.
\item{} Forklar hva som vil skje (og hvorfor) i reguleringssløyfen hvis FCE plutselig feiler med tilsvarende lav (4 mA) MV-signal.
\item{} Forklar hva som vil skje (og hvorfor) i reguleringssløyfen hvis FCE plutselig feiler med tilsvarende høy (20 mA) MV-signal.
\item{} Forklar hvordan du kan bekrefte tilstedeværelsen av et {\it åpent} brudd i en 4-20 mA signalkabel ved kun å bruke et voltmeter (ingen resistans- eller strømmåling tillatt!).
\item{} Forklar hvordan du kan bekrefte tilstedeværelsen av en {\it kortslutning} i en 4-20 mA signalkabel ved kun å bruke et voltmeter (ingen resistans- eller strømmåling tillatt!).  Hint: du må bryte kretsen.
\end{itemize}




\vfil \eject

\centerline{\bf En enkel lukket-sløyfe PID-tuningsprosedyre}

\vskip 10pt

Å tune en PID-regulator er på mange måter en kunst, og kan være ganske krevende for en nybegynner.  Det som følger er en primitiv (og forenklet for noen situasjoner!) prosedyre du kan bruke på mange prosesser.

\vskip 10pt

\noindent
{\bf Steg 1}

\underbar{Forstå prosessen du prøver å regulere.}  Hvis du ikke har en grunnleggende forståelse av prosessens natur, er det meningsløst -- til og med farlig -- å endre regulatorinnstillinger.  Her er en enkel sjekkliste før du rører regulatoren:

\begin{itemize}
\item{} Hva er prosessvariabelen, og hvordan måles den?
\item{} Hva er det endelige reguleringsorganet, og hvordan påvirker det prosessvariabelen?
\item{} Hvilke sikkerhetsfarer finnes i denne prosessen knyttet til regulering (f.eks. eksplosjonsfare, størkning, produksjon av farlige biprodukter, osv.)?
\item{} Hvilket reguleringsopplegg brukes?  Er dette en enkel, frittstående PID-algoritme, eller er den koblet til andre algoritmer eller spesialfunksjonsblokker?  Hvis dette ikke er en enkel algoritme, sørg for at du vet hvordan den skal reagere under alle tenkelige prosessforhold.
\item{} Hvor mye kan jeg "bump"-e prosessen mens jeg tuner regulatoren og observerer responsen?
\item{} Hvordan settes regulatoren i "manuell" modus, i tilfelle jeg må ta over kontrollen?
\item{} Hvis regulatoren forårsaker en farlig situasjon, hvordan stenger du prosessen ned?
\end{itemize}

\vskip 10pt

\noindent
{\bf Steg 2}

\underbar{Forstå hva innstillingene på regulatoren gjør.}  Er regulatoren konfigurert for gain eller proportional band?  Minutter per repetisjon eller repetisjoner per minutt?  Bruker den reset windup-grenser?  Reagerer derivasjon på feil eller bare PV?  Du må forstå hvordan PID-verdiene påvirker regulatorens virkemåte hvis du skal avgjøre hvilken vei (og hvor mye) du skal justere dem!  I dagene med analoge elektroniske og pneumatiske regulatorer anbefalte jeg teknikere å tegne små piler ved hvert justeringshjul som viste hvilken vei man skulle vri for mer aggressiv handling -- slik slapp de å blande gain vs PB, rep/min vs min/rep, osv.: alt de trengte å tenke på var "mer" eller "mindre" av hver handling.

\vskip 10pt

\noindent
{\bf Steg 3}

\underbar{"Bump" den manipulerte variabelen (endelig reguleringsorgan) manuelt for å lære hvordan prosessen responderer.}  Akkurat nå er {\it du} regulatoren!  Det du må gjøre er å justere prosessen for å lære hvordan den responderer: er det en integrerende prosess, en selvregulerende prosess eller en runaway-prosess?  Er det betydelig dødgang eller hysterese?  Er responsen lineær og konsistent?  Mange prosesskontrollproblemer skyldes faktorer som ikke er regulatoren, og dette "manuelle test"-steget er en viktig diagnostisk teknikk for å vurdere disse faktorene.

\vskip 10pt

\noindent
{\bf Steg 4}

\underbar{Sett PID-konstanter til "minimale" innstillinger og bytt til automatisk modus.}  Dette betyr gain mindre enn 1, ingen integrerende handling (0 rep/min eller maksimal min/rep), ingen derivat-handling, og ingen filtrering.

\vskip 10pt

\noindent
{\bf Steg 5}

\underbar{"Bump" settpunktet og se regulatorens respons.}  Dette tester regulatorens evne til å håndtere prosessen på egen hånd.  Det du ønsker er en respons som er rimelig rask uten for mye oversving eller undersving, og uten overdreven oscillasjon.  Prosessens natur og kravene til kvalitet vil avgjøre hva som er "for mye" responstid, over/undersving og svingning.

\vskip 10pt

\noindent
{\bf Steg 6}

\underbar{Øk eller reduser aggressiviteten i reguleringshandlingen basert på resultatene fra steg 5.}

\vskip 10pt

\noindent
{\bf Steg 7}

\underbar{Gjenta steg 5 og 6 for P, I og D, én om gangen, i den rekkefølgen.}  Med andre ord, still regulatoren først inn som en P-only regulator, legg så til integrerende handling (PI), deretter derivasjon (PID), etter behov.

\filbreak

\centerline{\bf Forbehold}

\vskip 10pt

Prosedyren beskrevet her er {\it svært} grov, og bør bare brukes som en students første forsøk på PID-tuning, på en trygg "demonstrasjons"-prosess.  Den antar at prosessen hovedsakelig responderer på proporsjonal (P-only) handling, noe som ikke nødvendigvis gjelder for alle prosesser.  Den gir heller ikke konkrete råd for tuning basert på resultatene i steg 3, noe som er kjennetegnet på en erfaren PID-tuner.  Med studier, øvelse og tid vil du lære hvilke prosess-typer som responderer best på P-, I- og D-handling, og du vil da kunne velge intelligente parametere å justere og hvilke lukket-sløyfe-oppførsel du skal se etter.

\vskip 10pt

\underbar{file i00027}
%(END_QUESTION)





%(BEGIN_ANSWER)


%(END_ANSWER)





%(BEGIN_NOTES)

\noindent
{\bf Sløyfediagram / inspeksjoner:}

Jeg anbefaler sterkt å krysse av studentenes sløyfediagrammer mens du inspiserer sløyfen deres (sjekker for sikker kabling, korrekt rørføring, god rør-/kanalinstallasjon, osv.) sammen med dem.  La alle teammedlemmer ta deg med på en "omvisning" av den ferdige sløyfen, der hvert medlem forklarer en annen del av sløyfen som du velger, mens de bruker sitt eget sløyfediagram som guide.  Mens en student forklarer sin del av sløyfen kan du kontrollere de andre studentenes sløyfediagrammer for nøyaktighet.  Dette sparer ikke bare tid ved å slå sammen oppgaver for sløyfeinspeksjon og sløyfediagram-verifikasjon, men det sikrer også at studentene faktisk kan knytte sløyfediagrammene til sløyfen de har bygget og formulere denne forståelsen for deg.

%INDEX% Lab exercise, building a complete process control loop (with networked controller)

%(END_NOTES)
